% !TEX root = ../main.tex
\chapter{Cơ sở lý thuyết và công nghệ sử dụng}

\section{Kiến trúc ứng dụng web toàn ngăn xếp}

Ứng dụng web toàn ngăn xếp gồm lớp giao diện trên trình duyệt, lớp dịch vụ xử lý nghiệp vụ và lớp lưu trữ. Trong RabbitCV, trình duyệt tải ứng dụng React, gửi yêu cầu HTTP đến REST API và thiết lập kết nối Socket.IO khi cần nhắn tin. Express xác thực, xử lý nghiệp vụ và làm việc với MongoDB thông qua Mongoose. Cách phân lớp này giúp giao diện và máy chủ phát triển tương đối độc lập, đồng thời cho phép kiểm tra quyền truy cập tại server.

\section{MERN Stack}

MERN là tổ hợp MongoDB, Express, React và Node.js. JavaScript được sử dụng xuyên suốt ở client và server, giúp giảm chi phí chuyển đổi ngôn ngữ và thuận tiện khi trao đổi dữ liệu JSON. Đối với một đồ án cá nhân, hệ sinh thái thống nhất giúp tập trung thời gian vào nghiệp vụ CV và tuyển dụng.

MERN không phải lựa chọn tốt nhất cho mọi hệ thống. Cơ sở dữ liệu quan hệ có lợi thế về ràng buộc tham chiếu và giao dịch; các backend có cấu trúc nghiêm ngặt có thể phù hợp hơn với nhóm lớn. RabbitCV chọn MERN vì dữ liệu CV có nhiều nhóm thuộc tính lồng nhau, giao diện có tính tương tác cao và phạm vi triển khai phù hợp với một ứng dụng JavaScript toàn ngăn xếp.

\section{React và quản lý trạng thái}

React xây dựng giao diện từ các component và cập nhật màn hình dựa trên state \cite{reactdocs}. Trong RabbitCV, các component dùng lại gồm header, CV preview, form CV, chat box, notification bell và các thành phần bảo vệ route. \texttt{AuthContext} là nguồn dữ liệu dùng chung cho người dùng hiện tại, trạng thái kiểm tra phiên và các thao tác đăng nhập, đăng xuất.

So với việc mỗi trang tự gọi API lấy thông tin người dùng, một context thống nhất giúp giảm trạng thái không nhất quán giữa header, route bảo vệ và trang chức năng. \texttt{SavedJobsContext} quản lý việc đã lưu; \texttt{SocketContext} quản lý kết nối Socket.IO, trạng thái tin nhắn chưa đọc và các thao tác đánh dấu đã đọc.

\section{Vite và React Router}

Vite cung cấp máy chủ phát triển và quá trình build frontend nhanh, hỗ trợ cập nhật nóng trong quá trình lập trình \cite{vitedocs}. RabbitCV sử dụng Vite để chạy frontend tại cổng 4000 và proxy các đường dẫn \texttt{/api} và \texttt{/socket.io} sang backend tại cổng 5000. Nhờ đó, frontend gọi đường dẫn tương đối và giảm sai khác giữa môi trường phát triển và triển khai.

React Router đảm nhiệm điều hướng phía client. Các route được chia theo nhóm công khai, ứng viên, doanh nghiệp và admin. \texttt{RoleAccessGuard} ngăn người dùng truy cập các khu vực không phù hợp với vai trò, nhưng việc kiểm tra quan trọng vẫn được thực hiện lại tại backend.

\section{Tailwind CSS và hệ thống giao diện}

Tailwind CSS cung cấp các lớp tiện ích để biểu diễn màu sắc, khoảng cách, lưới, kích thước và trạng thái tương tác \cite{tailwinddocs}. Cách tiếp cận này giúp RabbitCV giữ sự thống nhất giữa trang ứng viên, cổng doanh nghiệp và khu vực quản trị.

CSS riêng vẫn được dùng cho những phần có yêu cầu đặc thù như vùng in CV A4, trình xem PDF, thanh cuộn, kích thước responsive và các component rich text. Tiptap hỗ trợ trình soạn thảo nội dung có định dạng; Heroicons cung cấp biểu tượng cho các thao tác giao diện.

\section{Node.js và Express}

Node.js sử dụng mô hình I/O bất đồng bộ, phù hợp với máy chủ có nhiều thao tác chờ như truy vấn database, gửi phản hồi và giao tiếp Socket.IO. Express cung cấp middleware, định tuyến và cách tổ chức REST API gọn \cite{expressdocs}.

Trong phiên bản hiện tại, \texttt{server/server.js} vẫn là entry point lớn và chứa phần lớn route nghiệp vụ. Một số phần đã được tách thành các nhóm route, service, utility và cron. Đây là cách tổ chức phù hợp với phạm vi đồ án, nhưng việc tách controller và service theo domain là hướng cải thiện cần thực hiện khi hệ thống mở rộng.

\section{MongoDB và Mongoose}

MongoDB lưu tài liệu BSON và hỗ trợ cấu trúc lồng nhau \cite{mongodbdocs}. Điều này phù hợp với CV vì một Resume có các nhóm học vấn, kinh nghiệm, kỹ năng, dự án, chứng chỉ, ngôn ngữ và sở thích. Mongoose cung cấp schema, kiểu dữ liệu, giá trị mặc định, validation và index \cite{mongooseDocs}.

RabbitCV sử dụng ObjectId để liên kết người dùng, CV, công việc và hồ sơ ứng tuyển. Một số index quan trọng gồm cặp \texttt{jobId}--\texttt{candidateId} duy nhất trong Application để ngăn ứng tuyển trùng, cặp \texttt{recipient}--\texttt{isRead}--\texttt{createdAt} trong Notification để tối ưu truy vấn thông báo và cặp \texttt{userId}--\texttt{room} trong ChatReadState hoặc HiddenChat để lưu trạng thái theo người dùng.

\section{REST API và xác thực JWT}

REST tổ chức chức năng quanh tài nguyên và dùng phương thức HTTP biểu đạt thao tác. RabbitCV dùng \texttt{GET} để đọc, \texttt{POST} để tạo hoặc thực hiện hành động, \texttt{PUT} hoặc \texttt{PATCH} để cập nhật và \texttt{DELETE} để xóa hoặc bỏ liên kết. Phản hồi JSON giúp frontend xử lý thống nhất.

JWT được tạo sau khi đăng nhập thành công và lưu trong cookie HTTP-only. JavaScript phía trình duyệt không đọc trực tiếp token, giảm một phần rủi ro lộ token qua mã giao diện. Backend dùng middleware để kiểm tra token, người dùng hiện tại, vai trò và quyền sở hữu trước khi thực hiện thao tác nhạy cảm \cite{jwtdocs}.

Ngoài xác thực, server còn sử dụng bcryptjs để băm mật khẩu, CORS để kiểm soát nguồn truy cập, Helmet để thiết lập header bảo mật và express-rate-limit để giới hạn các request chung, đăng nhập, đăng ký và nhóm API AI.

\section{Socket.IO và tác vụ theo lịch}

Socket.IO cung cấp giao tiếp hai chiều dựa trên sự kiện, room và cơ chế kết nối lại \cite{socketiodocs}. Middleware Socket.IO đọc cookie, xác minh JWT và đưa socket vào room riêng theo người dùng. Khi có tin nhắn mới, server lưu Message vào MongoDB và phát sự kiện \texttt{newMessage} đến bên nhận.

Thông báo chưa đọc được truy vấn định kỳ vì dữ liệu nhỏ và không yêu cầu độ trễ tuyệt đối. Ngược lại, chat sử dụng Socket.IO vì cần phản hồi gần thời gian thực. \texttt{node-cron} chạy tác vụ kiểm tra tin sắp hết hạn hoặc đã hết hạn, cập nhật trạng thái Job và tạo Notification tương ứng.

\section{PDF, PDF.js và khổ giấy A4}

CV tạo trên web được trình bày bằng HTML/CSS rồi in qua trình duyệt. Quy tắc \texttt{@page}, kích thước 210 x 297 mm, lề trang và các thuộc tính in được dùng để kiểm soát đầu ra.

CV tải lên phải là PDF, có data URL hợp lệ, có header PDF và dung lượng không vượt quá 5 MB. Hệ thống từ chối DOC, DOCX, ảnh và nội dung giả PDF. PDF được lưu trong trường \texttt{pdfData} dưới dạng base64 trong phạm vi đồ án; danh sách CV không trả trường này để giảm băng thông.

Ở giao diện doanh nghiệp, PDF.js kết xuất từng trang vào vùng hiển thị riêng thay vì dùng iframe hoặc thanh công cụ PDF gốc \cite{pdfjsdocs}. Cách làm này giúp giao diện xem CV nhất quán, nhưng làm tăng chi phí xử lý phía trình duyệt và chưa giải quyết hoàn toàn bài toán lưu trữ tệp ở quy mô lớn.

\section{Tích hợp AI có kiểm soát}

Các chức năng AI được gọi từ Express backend thông qua API tương thích OpenAI của Groq. Frontend không nhận khóa API. Các luồng hiện có gồm viết, cải thiện hoặc rút gọn tiểu sử; review CV; điền nội dung CV; chấm mức độ phù hợp giữa CV và tin tuyển dụng; và luyện phỏng vấn theo từng lượt.

Trước khi gửi, dữ liệu CV được làm sạch và loại thông tin liên hệ trực tiếp. Backend áp dụng quota theo ngày, rate limit, timeout, schema đầu ra có cấu trúc và bộ lọc kết quả. Model \texttt{AiUsage} chỉ lưu thống kê như feature, model, trạng thái, hash đầu vào, số token, độ trễ và mã lỗi; không lưu prompt hoặc toàn bộ nội dung CV.

AI chỉ mang tính hỗ trợ tham khảo. Kết quả review, điểm phù hợp và điểm phỏng vấn không được xem là quyết định tuyển dụng thực tế. Người dùng phải kiểm tra và chủ động áp dụng nội dung trước khi lưu vào CV.

\section{Công cụ phát triển và kiểm thử}

Git được dùng để quản lý phiên bản và branch. ESLint kiểm tra quy tắc mã nguồn. Node Test Runner thực hiện kiểm thử tự động cho các hàm xử lý ngày tháng, PDF, dữ liệu AI, trạng thái việc làm, rich text, thông báo và các luồng giao diện quan trọng. Vite build frontend để phát hiện lỗi tích hợp tĩnh trước khi triển khai.
